\section{Section 3: The conditions on the eigenvalues}

\begin{lemma}
  \label{lem:s3_1}
  \lean{TwoControl.section3_lemma_3_1}
  \leanok
  Suppose $u_0, u_1$ are complex numbers with $|u_0| = |u_1| = 1$.
  An $8 \times 8$ unitary $U$ commutes with
  $\operatorname{Diag}(u_0, u_1) \otimes I_2 \otimes I_2$
  if and only if either $u_0 = u_1$ or
  $U = |0\rangle\langle 0| \otimes V_{00} + |1\rangle\langle 1| \otimes V_{11}$
  for some $4 \times 4$ unitaries $V_{00}, V_{11}$.
\end{lemma}

\begin{proof}
  \leanok
  Decompose $U$ into a $2 \times 2$ block matrix of $4 \times 4$ submatrices
  $U = |0\rangle\langle 0| \otimes V_{00}
     + |0\rangle\langle 1| \otimes V_{01}
     + |1\rangle\langle 0| \otimes V_{10}
     + |1\rangle\langle 1| \otimes V_{11}$.
  Similarly decompose $\operatorname{Diag}(u_0, u_1) \otimes I_2 \otimes I_2$
  as a block-diagonal matrix with blocks $u_0 \cdot I_4$ and $u_1 \cdot I_4$.
  Compute $UW$ and $WU$ and compare blocks:
  commutativity forces $u_1 \cdot V_{01} = u_0 \cdot V_{01}$
  and $u_0 \cdot V_{10} = u_1 \cdot V_{10}$.
  When $u_0 \neq u_1$, scalar cancellation implies $V_{01} = V_{10} = 0$,
  giving the block-diagonal form.
  Unitarity of $U$ then forces $V_{00}$ and $V_{11}$ to be individually unitary.
  The converse is straightforward.
\end{proof}

\begin{lemma}
  \label{lem:s3_2}
  \lean{TwoControl.section3_lemma_3_2}
  \leanok
  Suppose $u_0, u_1$ are complex numbers with $|u_0| = |u_1| = 1$.
  There exist $2 \times 2$ unitaries $P, Q$ such that the eigenvalues
  of $P \otimes Q$ are $\{1, 1, u_0, u_1\}$ (as a multiset, up to permutation)
  if and only if either $u_0 = u_1$ or $u_0 u_1 = 1$.
\end{lemma}

\begin{proof}
  \leanok
  (\emph{Forward direction.})
  Diagonalize $P$ and $Q$ via the spectral theorem
  to write $P = V_P D_P V_P^\dagger$ and $Q = V_Q D_Q V_Q^\dagger$.
  The eigenvalues of $P \otimes Q$ are the pairwise products
  of eigenvalues of $P$ and $Q$.
  These four products must be a permutation of $\{1, 1, u_0, u_1\}$.
  Exhaustive case analysis over all $24$ permutations of four elements
  shows each case yields either $u_0 = u_1$ or $u_0 u_1 = 1$.

  (\emph{Backward direction.})
  If $u_0 = u_1$, take $P = \operatorname{Diag}(1, u_1)$ and $Q = I_2$.
  If $u_0 u_1 = 1$, take $P = \operatorname{Diag}(1, u_0)$
  and $Q = \operatorname{Diag}(1, u_1)$
  with a suitable diagonalizing unitary.
\end{proof}

\begin{lemma}
  \label{lem:s3_3}
  \lean{TwoControl.section3_lemma_3_3}
  \leanok
  Suppose $u_0, u_1$ are complex numbers with $|u_0| = |u_1| = 1$.
  There exist a $2 \times 2$ unitary $P$ and complex numbers $c, d$
  such that the eigenvalues of
  $(I_2 \otimes P) \cdot C(\operatorname{Diag}(u_0, u_1))$
  are $\{c, c, d, d\}$ (each with multiplicity two)
  if and only if either $u_0 = u_1$ or $u_0 u_1 = 1$.
  Here $C(\operatorname{Diag}(u_0, u_1))$ denotes the
  controlled-$\operatorname{Diag}(u_0, u_1)$ gate.
\end{lemma}

\begin{proof}
  \leanok
  (\emph{Forward direction.})
  Rewrite $(I_2 \otimes P) \cdot C(\operatorname{Diag}(u_0, u_1))$
  as the direct sum $P \oplus (P \cdot \operatorname{Diag}(u_0, u_1))$.
  Diagonalize $P$ and $P \cdot \operatorname{Diag}(u_0, u_1)$
  separately via the spectral theorem.
  The eigenvalues of the direct sum are related by a permutation
  to those of $\operatorname{Diag}(c, d, c, d)$.
  Exhaustive case analysis on all $24$ permutations of four eigenvalues:
  if both eigenvalues of $P$ equal the same value $c$,
  then $P$ is a scalar multiple of the identity,
  and a cancellation argument forces $u_0 = u_1$.
  In the remaining cases, a determinant argument
  using $\det(AB) = \det(A) \cdot \det(B)$ yields $u_0 u_1 = 1$.

  (\emph{Backward direction.})
  If $u_0 = u_1$, take $P = I_2$.
  If $u_0 u_1 = 1$, take $P = \operatorname{Diag}(1, u_0)$.
\end{proof}